

Compression is domain-dependent: coding agents require structural fidelity and state continuity; web and GUI agents must preserve heterogeneous multimodal observations; DeepResearch agents depend on deferred evidence recovery across long horizons. Figure~\ref{fig:domain_specific_regimes} summarizes these regimes.

% The practical value of a compression strategy is domain-dependent. Coding agents require high structural fidelity and explicit state continuity; web and GUI agents face heterogeneous multimodal observations; research agents depend most on deferred evidence recovery across long horizons. Figure~\ref{fig:domain_specific_regimes} summarizes these regimes.


\noindent\textbf{Coding Agents}
Coding places the strongest pressure on structure preservation, exact evidence retention, and reliable recovery. Accordingly, methods such as SWE-Pruner~\citep{wang2026sweprunerselfadaptivecontextpruning}, LongCodeZip~\citep{shi2025longcodezip}, and CAT~\citep{liu2025contexttoolcontextmanagement} emphasize structure-aware pruning and explicit state retention.
Table~\ref{tab:coding_agents_pipeline} compares recent production-grade coding agents through the unified operational pipeline, using only public evidence and high-confidence product signals.
The contrast is instructive. Claude Code emphasizes \textbf{Compress} through reactive compaction, Cursor places more weight on \textbf{Store} and \textbf{Recover} through repository indexing and retrieval, and Codex appears most grounded in \textbf{Select} plus workspace continuation. Across all three, the common pattern is not aggressive lossy shortening, but coordinated pipeline control under tight context budgets.


\noindent\textbf{Web \& GUI Agents}
Web and GUI agents operate over heterogeneous, partially redundant observations. Methods such as LCoW~\citep{lee2025learningcontextualizewebpages} and PAL-UI~\citep{liu2025paluiplanningactivelookback} exploit redundancy, but representation-level compression must still preserve modality and interface cues.

\noindent\textbf{DeepResearch  Agents}
Research and deep-search agents are defined less by structural fidelity than by delayed evidence use. Their main risk is that abstraction hides evidence that matters only later. ReSum~\citep{wu2026resumunlockinglonghorizonsearch} and COMPASS~\citep{wan2025compassenhancingagentlonghorizon} favor summarization, whereas Step-DeepResearch~\citep{hu2025stepdeepresearchtechnicalreport} and WebWeaver~\citep{li2025webweaverstructuringwebscaleevidence} externalize or retrieve evidence to preserve long-horizon dependencies.
